\section{Introduction}
\label{sec:intro}

Graph neural networks (GNNs) have become the standard tool for learning on relational data, achieving strong results on node classification, link prediction, and recommendation~[GCN, GraphSAGE, GAT]. Their effectiveness rests on a quiet assumption: that every node carries an informative feature vector, which message passing then refines by aggregating information across edges. In practice this assumption often fails. Nodes join a network before a profile is filled in, newly listed items lack descriptions, sensors drop out, and privacy constraints withhold attributes. The result is a graph in which a substantial fraction of nodes have \emph{no observed features at all}. Under such \textbf{attribute-missing} conditions, message passing degrades sharply: a node whose neighborhood is largely uninformative has little to aggregate, and the deficiency compounds with each additional missing hop. The problem is acute at scale: on a citation graph such as ogbn-arxiv (169K nodes), when a large fraction of nodes, say 80\%, have no observed features, most local neighborhoods are almost entirely empty, and a standard GNN has little signal left to propagate.

We study this setting through the lens of the \emph{downstream task} rather than reconstruction. Much prior work treats missing attributes as an imputation problem (fill in the missing features, then run a standard model) and measures success by how faithfully the raw features are recovered~[SAT, SVGA]. Yet accurate reconstruction is neither necessary nor sufficient for an accurate decision boundary. What ultimately matters is whether the learned representations remain discriminative when features are scarce. We therefore adopt \textbf{learning under missing node features}: given a graph where whole feature vectors are absent, learn node representations that support accurate downstream node classification.

Approaches to missing node features fall into a few broad families, each with a characteristic limitation. \emph{Imputation-based} methods reconstruct features before applying a GNN; propagation variants~[FP] are simple and scalable, while learned generative imputers~[SAT, SVGA] are more expressive but rely on dense or per-node operations that do not scale. \emph{Imputation-free} methods modify the aggregation to tolerate missing entries without filling them~[PaGCN, GCNMF]. \emph{Per-node parameterization} assigns each missing node its own learnable embedding~[MATE], which risks memorization when no supervision reaches those parameters. 
Most recently, \emph{self-supervised and contrastive} methods learn missingness-robust representations from pretext tasks instead of relying on reconstruction quality~[AmGCL, MOBA]. But these methods are applied largely by analogy: an auxiliary objective is borrowed from the graph-SSL literature~[GraphMAE, BarlowTwins, MVGRL] and added in the hope that it helps. What is missing is a principled account of \emph{which} self-supervised objectives help a GNN learn under missing features, and \emph{when}.

This question is a natural continuation of our own prior work~[NCSSL], which introduced a feature-based neighborhood self-supervision for graph attribute \emph{imputation}, where it helped precisely because the goal was feature reconstruction and the objective supervised the same feature space. Here we ask a distinct question: does that auxiliary signal also improve the \emph{downstream classification} boundary, and under what conditions? And how well they apply to bigger emebedding datatsets? Since reconstruction and classification are different objectives, the answer is not implied by the earlier result. We therefore treat self-supervision as an object of study rather than an assumption extending the earlier framework.

We introduce \textbf{NESS}, a GNN for learning under missing node features that couples (i)~a Feature-Propagation prefill that gives every missing node long-range structural signal before training, and (ii)~a set of \textbf{self-supervised objectives} layered on a well-configured encoder. NESS outperforms imputation-based, imputation-free, and per-node-parameter baselines on downstream classification, with the advantage most pronounced under severe missingness and at a scale where generative imputation is infeasible.

Beyond the method, our central contribution is a systematic study of self-supervised objectives for learning under missingness. We find that self-supervision generally helps, but the size of its benefit depends on the setting. Two factors modulate it. First, the severity of missingness: self-supervision contributes most when the encoder is information-starved, and its advantage grows as features become scarcer. Second, the nature of the features: neighborhood-centric objectives are more effective on dense, continuous features, but less on sparse binary bag-of-words attributes. We further find that the benefit varies with the objective itself, and analyze which properties make an objective effective (Section~\ref{sec:exp-why}). A label-aware objectives is the most reliable objective, and it transfers as an enhancement to other backbones. This also motivates our focus on large dense-embedding graphs, the regime prior small-scale binary benchmarks least represent.

Our contributions are as follows:
\begin{itemize}
    \item We introduce NESS, a graph neural network for learning under missing node features that couples a Feature-Propagation prefill with pluggable self-supervised objectives on a well-configured encoder, outperforming imputation-based, imputation-free, and per-node-parameter baselines on downstream node classification, with the largest advantage under severe missingness and at a scale where generative imputation is infeasible.
    \item We provide a empirical study of when and why self-supervision helps under missingness, showing that its benefit generally holds but is modulated by the severity of missingness, the type of features, and the characteristics of the objective, which we analyze rather than assume.
    \item We identify two reliable objectives for our setting: a label-aware neighbor-histogram and a structure based objective. We also show that it transfers as an enhancement to other GNN backbones.
\end{itemize}


In the following sections we review related work on graph learning under weak information (Section~\ref{sec:related}), discuss the different self-supervised objectives used in graph learning (Section~\ref{sec:ssl-objectives}), introduce our method (Section~\ref{sec:method}), and conduct comprehensive experiments to address our research questions (Section~\ref{sec:experiments}).
